% !TEX root = main.tex
\section{Homomorphisms Preserves Measure Between Source Value Set}~\label{sec:homomorphisms-preserves-measure-between-source-value-set}

Let us return to an examination of the measure or density of $\mathbb{F}_{n}(b)$.

\begin{definition}
Denote
\[\mathbb{F}(b) = \{ a \mid a\rightsquigarrow b \}\cup\{b\} = \bigcup_{n = 0}^{\infty}\mathbb{F}_{n}(b)\]
as all numbers that will be convergent to $b$.
\end{definition}
\begin{proposition}~\label{prop:source-set-disjoint}
	If $a \not\rightsquigarrow b$ and $b \not\rightsquigarrow a$, $\mathbb{F}(a) \cap \mathbb{F}(b) = \varnothing$.
\end{proposition}
\begin{proposition}
	If $a\rightsquigarrow b$, $\mathbb{F}(a) \subset \mathbb{F}(b)$.
\end{proposition}
\begin{theorem}~\label{thm:infinite-layers}
	For $b\not\in\mathbb{S}$, we have $\forall N > 0$, $\mathbb{F}(b) - \bigcup_{n=0}^{N}\mathbb{F}_{n}(b) \ne \varnothing$.
\end{theorem}
\begin{proof}
	Divide it into two situations:\newline
	\begin{itemize}
		\item[If $b \not\rightsquigarrow b$:] 
		Since $b\not\in\mathbb{S}$, from \cref{prop:all-source-value-set-has-primitive-source}, we have $\mathbb{F}_{N}(b) \not\subset \mathbb{S}$. That means $\mathbb{F}_{N + 1}(b) \ne \varnothing$.
		From \cref{prop:non-cycle-layer-independent}, $\forall N \ge 0$ there is
		\[\mathbb{F}(b) - \bigcup_{n=0}^{N}\mathbb{F}_{n}(b) \supset \mathbb{F}_{N+1}(b) - \bigcup_{n=0}^{N}\mathbb{F}_{n}(b) = \mathbb{F}_{N+1}(b) \ne \varnothing.\]
		\item[If $b \rightsquigarrow b$:] Suppose it's a $t$-period cycle, then $b \in\mathbb{F}_{t}(b)$. From \cref{thm:layers-of-reduce-space-and-source-value-set}, we can just pick $a = \beta^{\Phi(1)}\cdot b + \frac{\gamma}{\alpha}\cdot(\beta^{\Phi(1)} - 1)\in\mathbb{F}_{t}(b)$ that $b\not\rightsquigarrow a$ but $a \rightsquigarrow b$. From previous discussion, $\mathbb{F}(a)$ can't be covered within finity layers, thus $\mathbb{F}(b) \supset \mathbb{F}(a)$ also can't be covered within finity layers.
	\end{itemize}
\end{proof}
% \begin{corollary}~\label{cor:card-of-source-value-set}
%     $\mathrm{card}(\mathbb{F}_{0}(B)) = \mathrm{card}(B)$, $\mathrm{card}(\mathbb{F}_{n+1}(B)) = {\aleph_{0}}\times\mathrm{card}(\mathbb{F}_{n}(B))$.
% \end{corollary}

\cref{thm:infinite-layers} demonstrates that we can never resolve the Collatz conjecture through any method of case analysis, even if we partition the problem into various infinite sets.
Hence, we shall instead attempt to identify certain homomorphisms in the hope of capturing specific measure-theoretic properties that, we trust, may ultimately aid us in unraveling this conjecture.

From Tao's result\cite{tao2022orbitscollatzmapattain}, when $(\alpha, \beta, \gamma) = (3, 2, 1)$, we have the logarithm density of $\mathbb{F}(1)$ is $1$.
Suppose there's another number $b$ that $b \not\rightsquigarrow 1$, from \cref{prop:source-set-disjoint} we have $\mathbb{F}(b)\cap\mathbb{F}(1) = \varnothing$, that means the logarithm density of $\mathbb{F}(b)$ will be $0$.
Since $\mathbb{F}(b)$ can't be covered by finity layers.
Thus, if there a homomorphism that keeps the logarithm measure, or keeps the logarithm measure by layer count, then we can prove that $b$ will not exists. That means there's no one that will not fall into cycle can escape from $1$.

Thus, if we can prove there is a homomorphism that keeps the logarithm measure between those source value set, we proved \cref{conj:source-value-of-1-cover-all-odd-numbers}, which is equivalent to the classic Collatz conjecture.
Fortunately, from \cref{sec:varies-in-source-value-with-respect-to-the-reduce-sequence}, we can build such a homomorphism.

TODO